\chapter{Kết luận và hướng phát triển}
\label{ch:ket-luan}

Chương này tổng hợp kết quả theo bốn câu hỏi nghiên cứu, nêu đóng góp và bài học thiết kế, chỉ
rõ giới hạn của phép đo, và đề xuất hướng phát triển. Các kết luận được viết ở mức mà dữ liệu
cho phép: \emph{sức mạnh thi đấu trong cấu hình đã đo}, không phải \emph{tối ưu đã chứng minh}.

\section{Trả lời các câu hỏi nghiên cứu}
\label{sec:kl-rq}

\paragraph{RQ1. Trong cùng thời gian mỗi nước, MCTS có thắng nhiều hơn Alpha-Beta không?}
Trong các phép đo hiện có, Alpha-Beta mạnh hơn MCTS rất rõ. Ở ngân sách cố định, MCTS rollout
ngẫu nhiên thua toàn bộ 200 ván trước L2 và L3 ở 100, 400 và 1\,600 mô phỏng; chỉ ở 6\,400 mô
phỏng mới thắng 7,5\% trước L2 và 2,0\% trước L3. Trong RR8, L4 MCTS 400 mô phỏng đạt 14,3\%
tổng và thua L2 (43,6\%). Rollout nhẹ không cải thiện điểm, chỉ làm chậm khoảng 4 lần. Phân
tích lỗi cho thấy khoảng 90\% ván thua của MCTS là do không ngăn đe doạ kép sớm, một cấu trúc
cần cây sâu khoảng 4 ply trong khi ngân sách bị chia cho hàng chục nước gốc.

Tuy nhiên, báo cáo \emph{chưa} chạy chế độ cùng thời gian mỗi nước cho đủ tám tác tử. Đây là
khoảng trống quan trọng: Alpha-Beta có hàm lượng giá viết tay và nước bắt buộc nên mỗi nút đắt
hơn, còn MCTS có thể dùng nhiều mô phỏng hơn trong cùng 1 giây. Kết luận hiện tại chỉ đúng cho
ngân sách cố định theo loại thuật toán. Để trả lời đầy đủ RQ1, cần chạy thêm chế độ A với cùng
1\,000 ms/nước và báo cáo cả điểm lẫn thời gian.

\paragraph{RQ2. Self-play có cải thiện so với bắt chước chuyên gia không?}
Có, nhưng mức cải thiện phụ thuộc vào giai đoạn và đối thủ. Ba lớp bằng chứng:
\begin{itemize}
  \item \textbf{Cập nhật policy bằng RL:} RL v3 cải thiện khoảng 12--16 điểm phần trăm so với
        SL v3 trước L2/L3 trong phép đánh giá checkpoint (81,0\% so với 66,0\% trước L2; 37,0\%
        so với 21,0\% trước L3).
  \item \textbf{Thêm tìm kiếm khi chơi:} với cùng mạng L5, PUCT 200 mô phỏng nâng điểm trước L3
        từ 21,0\% lên 58,5\%, tức 37,5 điểm phần trăm. PUCT 800 mô phỏng đạt 67,0\%.
  \item \textbf{Huấn luyện policy--value bằng tự đấu có tìm kiếm:} L8 vòng 37 thắng L7 82,5\%
        trên 20 ván và thắng L3 68,75\% trên 40 ván.
  \item \textbf{Giải vòng tròn đầy đủ:} L6 đạt 61,9\% so với L5 50,2\%; L7 đạt 82,0\% và L8
        đạt 89,8\% so với L3 58,3\%. Đây là bằng chứng mạnh nhất rằng tìm kiếm có mạng vượt
        Alpha-Beta viết tay trong cấu hình này.
\end{itemize}
Ba phép đo tách ba mức đóng góp khác nhau, không nên cộng chúng thành một con số duy nhất. Hạn
chế là mỗi phép đo dùng một seed, cỡ mẫu cuối nhỏ (20--40 ván), và tập khai cuộc đánh giá cuối
từng được dùng trong gating. Chưa có lặp nhiều seed để định lượng độ ổn định của từng thay đổi.

\paragraph{RQ3. Cắt tỉa alpha-beta và sắp xếp nước giảm chi phí bao nhiêu mà vẫn giữ nước đi?}
A1 cho thấy cắt tỉa giảm khoảng 20,4 lần số nút và 20,9 lần thời gian, giữ nguyên giá trị gốc
trên 100\% số thế. Đây là bằng chứng mạnh nhất về tính đúng đắn của cài đặt. A2 cho thấy sắp
xếp nước giảm khoảng 24,5 lần số nút trong cùng 1 giây, nhưng không giảm thời gian vì ngân sách
thời gian cố định; nó giúp tìm sâu hơn. A7 cho thấy bỏ thành phần phòng thủ thay đổi giá trị
gốc trên 90\% số thế, tức thành phần này không thể bỏ qua.

Hai hạn chế. Thứ nhất, benchmark A2 không ghi độ sâu đạt được, nên không tách được ``nhiều nút
hơn vì tìm sâu hơn'' với ``nhiều nút hơn vì TT kém hiệu quả''. Thứ hai, báo cáo chưa đo tỷ lệ
chọn cùng nước khi có nhiều nước đồng hạng; trường nhất quán của benchmark chỉ so giá trị gốc.
A7b và A8b (nếu chạy xong) bổ sung phần đối đầu, nhưng vẫn cần một phép đo cố định độ sâu để
tách chi phí khỏi chất lượng.

\paragraph{RQ4. Lợi thế đi trước trên 11×11 lớn đến mức nào giữa các tác tử cùng cấp?}
Lợi thế đi trước không phải một hằng số. Kết quả E3:
\begin{itemize}
  \item L1 Random: X thắng 52\%, O thắng 48\% — gần cân bằng, đúng với tác tử không biết khai
        thác thế mở đầu.
  \item L2 Rule: X thắng 75\%, O thắng 25\% — lệch mạnh.
  \item L3 Alpha-Beta: X thắng 66\%, O thắng 18\%, hoà 16\% — lệch mạnh nhưng có hoà.
  \item L4 MCTS: X thắng 50\%, O thắng 50\% — cân bằng tuyệt đối.
  \item L5 Policy SL: X thắng 62\%, O thắng 36\%, hoà 2\%.
  \item L6 Policy RL: X thắng 50\%, O thắng 36\%, hoà 14\%.
  \item L7 PUCT: X thắng 80\%, O thắng 14\%, hoà 6\% — lệch mạnh nhất.
  \item L8 AlphaZero: X thắng 50\%, O thắng 22\%, hoà 28\% — lợi thế trung bình nhưng nhiều hoà.
\end{itemize}
Kết quả này phản bác giả định ``lợi thế đi trước là hằng số của trò chơi''. Nó là thuộc tính
của \emph{cặp tác tử}: hai tác tử yếu có thể cân bằng, hai tác tử mạnh có thể lệch, và một tác
tử có thể lệch theo hướng này nhưng không lệch theo hướng khác. Trong RR8, mọi tác tử đều đạt
điểm cao hơn khi cầm X, nhưng mức chênh khác nhau: L8 95,0\% so với 84,6\%, L7 90,3\% so với
73,7\%, L3 68,9\% so với 47,7\%, L5 56,1\% so với 44,3\%. Trong các trận L8--L3, L8 cầm X đạt
92,5\% và cầm O 45,0\%. Trong 10\,000 ván tự đấu từ bàn trống của L8, X thắng 77,0\%, O thắng
9,2\%, hoà 13,8\%; nhưng khi tự đấu từ khai cuộc 2--4 nước (E3), L8 chỉ còn 50/22/28. Điều này
cho thấy khai cuộc là một biến kiểm soát mạnh: nó làm giảm lợi thế đi trước và làm kết quả phản
ánh sức mạnh tác tử tốt hơn. Kết luận chỉ áp dụng cho bộ khai cuộc và ngân sách đã dùng.

\section{Đóng góp}
\label{sec:kl-dong-gop}
Báo cáo có sáu đóng góp, tương ứng với sáu mục ở Chương~\ref{ch:gioi-thieu}:
\begin{enumerate}
  \item Một hệ thống chạy được từ đầu đến cuối, dùng chung engine, giao diện tác tử, arena,
        huấn luyện, log và giao diện web.
  \item Tám cấp tác tử có thể so sánh trực tiếp, từ Random đến PUCT/AlphaZero-inspired.
  \item Một giao thức đánh giá có kiểm soát nhiễu, với khoảng Wilson, bootstrap theo cặp khai
        cuộc và phân tích theo màu.
  \item Phát hiện và sửa lỗi quy ước dấu của value head; đo lại calibration trước và sau khi
        sửa.
  \item Kết quả thực nghiệm mới về lợi thế đi trước (E3) cho thấy nó không đơn điệu theo sức
        mạnh tác tử.
  \item Phân tích ván có kiểm chứng bằng engine, đếm cơ hội thắng/chặn và phân loại lỗi trên
        hàng nghìn ván, thay vì mô tả định tính.
\end{enumerate}

\section{Bài học thiết kế}
\label{sec:kl-bai-hoc}
\paragraph{Tách quy ước dấu ngay từ đầu.} Lỗi value head là bài học đắt nhất của dự án. Nhãn
kết quả ván có hai quy ước tự nhiên: tuyệt đối (X thắng) và tương đối (người đi thắng). Nếu
không cố định quy ước ở một nơi và kiểm tra bằng test, lỗi sẽ lan sang RL và PUCT, nơi value
được dùng làm baseline hoặc hàm lượng giá. Bài học: định nghĩa quy ước trong một module, viết
test kiểm tra dấu, và đo calibration theo cả hai quy ước.

\paragraph{Đo chi phí cùng chất lượng.} Một tác tử có thể thắng nhờ được cấp nhiều thời gian
hơn, không nhờ thuật toán tốt hơn. Mọi so sánh trong báo cáo đều ghi ngân sách, backend và
thời gian mỗi nước. Khi ngân sách là thời gian, cần báo cáo cả số nút/mô phỏng đạt được; khi
ngân sách là số nút/mô phỏng, cần báo cáo cả thời gian. Không có một phép đo nào trả lời cả hai
câu hỏi.

\paragraph{Kiểm chứng chiến thuật bằng engine.} Phân tích ván định tính dễ bị thiên lệch bởi
ấn tượng. Dùng chính \texttt{winning\_squares} để đếm cơ hội thắng/chặn giúp phân biệt ``thua vì bỏ
chặn'' với ``thua vì đối thủ mạnh hơn ở thế cờ''. Cách làm này có thể mở rộng thành một bộ
kiểm tra tự động cho mọi ván.

\paragraph{Đa dạng dữ liệu quan trọng hơn số epoch.} SL v1 và v2 dùng cùng kiến trúc, cùng số
epoch, nhưng v2 có 297\,098 mẫu với temperature 8 thay vì 27\,720 mẫu với temperature 0,5. Top-1
tăng từ 44,0\% lên 65,3\%. SL v3 thêm 30 epoch chỉ tăng 2,5 điểm. Bài học: khi trần dữ liệu
đã đạt, thêm epoch không giúp; phải đổi nguồn nhãn hoặc tăng đa dạng.

\paragraph{Chọn checkpoint theo tiêu chí định trước.} Điểm gating của L8 dao động 43--62\% và
không so sánh trực tiếp giữa các vòng vì đối thủ đương nhiệm thay đổi. Một quy trình tốt cần
một tập đánh giá giữ riêng, cố định, và tiêu chí chọn checkpoint viết trước khi chạy. Báo cáo
dùng vòng 37 vì đó là lần duyệt cuối, nhưng không khẳng định đó là checkpoint tốt nhất tuyệt
đối.

\section{Giới hạn}
\label{sec:kl-gioi-han}
\paragraph{Giới hạn của phép đo.} Cỡ mẫu 100--200 ván mỗi cặp cho khoảng tin cậy rộng; chênh
lệch dưới 10 điểm phần trăm thường chưa phân biệt được. Các mốc đánh giá dùng khai cuộc khác
nhau, nên đường cong huấn luyện có nhiễu. Một seed gốc chưa đo được độ ổn định giữa các seed.
Các trận cuối L8--L7 (20 ván) và L8--L3 (40 ván) nhỏ, và tập khai cuộc từng dùng trong gating.

\paragraph{Giới hạn của mô hình.} Mạng chỉ có khoảng 0,48 triệu tham số, huấn luyện trên một
GPU laptop. L8 chạy 10\,000 ván tự đấu trong 12 giờ, nhỏ hơn nhiều so với AlphaZero gốc. Kết
quả không phản ánh năng lực của các hệ thống gốc, và không nên dùng để dự đoán chúng.

\paragraph{Giới hạn của luật và bàn cờ.} Bàn 11×11, luật tự do, không áp dụng luật chặn hai đầu
hay Renju. Kết quả không ngoại suy trực tiếp sang bàn 15×15 hoặc các biến thể khác. Tập ứng
viên bán kính 2 là cắt tỉa tiến, có thể loại nước tốt; tác tử không tìm kiếm trên toàn bộ cây.

\paragraph{Giới hạn của phân tích.} Phân tích chiến thuật chỉ đếm cơ hội thắng ngay và chặn bắt
buộc, không đánh giá được nước đi tối ưu dài hạn. Các nhận xét ván dựa trên luật thắng ngay và
đe doạ nước kế tiếp, không dựa trên phân tích sâu nhiều nước. Một nước không bỏ chặn vẫn có
thể là sai về chiến lược; báo cáo không kết luận ngược lại.

\section{Hướng phát triển}
\label{sec:kl-huong}
\paragraph{Đánh giá cùng thời gian.} Chạy chế độ A với 1\,000 ms/nước cho đủ tám tác tử, báo
cáo điểm và số nút/mô phỏng đạt được. Đây là phần còn thiếu quan trọng nhất để trả lời RQ1.

\paragraph{Nhiều seed.} Lặp mỗi cấu hình với 3--5 seed gốc, báo cáo trung bình và khoảng bất
định giữa các seed. Điều này đặc biệt quan trọng cho L8, nơi quá trình tự đấu có phương sai lớn.

\paragraph{Tập khai cuộc giữ riêng.} Tách một tập khai cuộc chưa dùng trong phát triển/gating,
đánh giá lại các cặp chính trên tập đó. Đây là cách duy nhất để tránh thiên lệch chọn mẫu.

\paragraph{Tăng cường chiến thuật.} Thêm lớp luật thắng/chặn trước khi gọi mạng, như các engine
Gomoku mạnh. Kết quả phân tích lỗi cho thấy L5/L6 bỏ 4--12 cơ hội thắng ngay trên 90 cơ hội;
một lớp luật đơn giản có thể sửa phần lớn mà không cần huấn luyện thêm.

\paragraph{Mạng lớn hơn và dữ liệu tốt hơn.} Tăng số kênh/khối, dùng nhiều dữ liệu chuyên gia
hơn, hoặc chưng cất tri thức từ Alpha-Beta độ sâu cao hơn. SL v3 đã gần trần của dữ liệu hiện
tại; đổi nguồn nhãn là cách tiếp theo.

\paragraph{Đánh giá nhiều biến thể luật.} Chạy cùng hệ thống trên luật chặn hai đầu, Renju, hoặc
bàn 15×15 để xem kết luận nào phụ thuộc vào luật. Đây là hướng mở rộng tự nhiên nhưng nằm
ngoài phạm vi đồ án.

\section{Kết luận}
\label{sec:kl-ket}
Đề tài đã xây dựng một hệ thống Game AI đa cấp độ cho Caro 11×11, hiện thực tám cấp tác tử trên
cùng một engine, và đánh giá chúng bằng một giao thức có kiểm soát nhiễu. Kết quả chính là:
tri thức chiến thuật viết tay vẫn rất mạnh ở ngân sách nhỏ; tìm kiếm có mạng cải thiện rõ so
với policy-only; self-play có cải thiện so với bắt chước nhưng cần nhiều ván; và lợi thế đi
trước phụ thuộc vào cặp tác tử, không phải hằng số của trò chơi.

Báo cáo cũng chỉ rõ những gì chưa làm được: chế độ cùng thời gian, nhiều seed, tập khai cuộc
giữ riêng, A6/A7b/A8b đầy đủ, và đánh giá trên biến thể luật. Những phần này được liệt kê ở
Phụ lục~\ref{app:cho-bo-sung} để người đọc biết chính xác phạm vi bằng chứng. Cách trình bày
này có chủ đích: một báo cáo nghiên cứu không chỉ nêu kết quả mà còn nêu điều kiện để kết quả
đó có giá trị.